% ════════════════════════════════════════════════════════════════════════════
% CHƯƠNG 3: KẾT QUẢ CÀI ĐẶT VÀ TRÌNH DIỄN
% ════════════════════════════════════════════════════════════════════════════
\chapter{KẾT QUẢ CÀI ĐẶT VÀ TRÌNH DIỄN}

\section{Mô tả cài đặt các chức năng cốt lõi}

\subsection{Tổng quan kiến trúc code}

Ứng dụng tổ chức code theo mô hình \textbf{Feature-First Architecture} với 3 tầng rõ ràng:

\begin{figure}[H]
\centering
\begin{tikzpicture}[
  layer/.style={draw, rounded corners=8pt, minimum width=13cm, minimum height=1.5cm,
    align=center, font=\small},
  >=Stealth
]
\node[layer, fill=blue!12] (ui) at (0, 0) {
  \textbf{Presentation Layer (UI)}\\
  Pages, Screens, Widgets --- Material Design 3\\
  \texttt{ConsumerWidget} / \texttt{ConsumerStatefulWidget}
};
\node[layer, fill=green!12] (state) at (0, -2.5) {
  \textbf{State Layer (Riverpod)}\\
  \texttt{StateProvider}, \texttt{StateNotifierProvider}\\
  Quản lý và phân phối state cho UI
};
\node[layer, fill=orange!12] (data) at (0, -5) {
  \textbf{Data Layer (Repository + Service)}\\
  \texttt{TransactionsRepository}, \texttt{BudgetsRepository}, \texttt{AccountsRepository}\\
  \texttt{AuthService}, \texttt{AiFinanceService}
};
\node[layer, fill=red!8] (ext) at (0, -7.5) {
  \textbf{External Services}\\
  Cloud Firestore, Firebase Auth, FCM, Gemini AI, Supabase, SePay
};

\draw[->, thick] (ui) -- node[right, font=\footnotesize] {\texttt{ref.watch()}} (state);
\draw[->, thick] (state) -- node[right, font=\footnotesize] {\texttt{ref.read()}} (data);
\draw[->, thick] (data) -- node[right, font=\footnotesize] {REST / SDK} (ext);
\end{tikzpicture}
\caption{Kiến trúc 4 tầng của G13 Money}
\label{fig:architecture-layers}
\end{figure}

\subsection{Chức năng 1: Xác thực người dùng (AuthService)}

\texttt{AuthService} là singleton class xử lý toàn bộ logic xác thực. Đoạn code tiêu biểu dưới đây minh họa quy trình đăng ký tài khoản mới:

\begin{lstlisting}[caption={AuthService.register() --- Đăng ký tài khoản với rollback}]
static Future<UserModel> register({
  required String fullName,
  required String phone,
  required String email,
  required String password,
}) async {
  User? createdUser;
  try {
    // B1: Tao tai khoan Firebase Auth
    final credential = await _auth
      .createUserWithEmailAndPassword(
        email: email.trim(),
        password: password,
      );
    createdUser = credential.user;

    // B2: Tao du lieu ban dau trong Firestore
    //   - User document
    //   - Settings profile + preferences
    //   - 11 danh muc mac dinh
    await _createInitialUserDocs(
      user: createdUser!,
      fullName: fullName,
      phone: phone,
    );

    // B3: Doc lai UserModel tu Firestore
    _currentUser =
      await _upsertAndReadUserModel(createdUser);
    return _currentUser!;

  } on FirebaseAuthException catch (e) {
    throw Exception(_authErrorMessage(e.code));
  } catch (e) {
    // ROLLBACK: Xoa user Auth neu Firestore fail
    if (createdUser != null) {
      try { await createdUser.delete(); }
      catch (_) {}
      await _auth.signOut();
    }
    throw Exception('Tao nguoi dung that bai');
  }
}
\end{lstlisting}

\textbf{Giải thích kỹ thuật:}
\begin{itemize}
  \item \textbf{Atomic registration:} Đăng ký gồm 3 bước (Auth → Firestore docs → Read back). Nếu bước 2 thất bại, bước 1 được rollback (xóa user Auth) để đảm bảo consistency.
  \item \textbf{Batch write:} \texttt{\_createInitialUserDocs()} sử dụng Firestore \texttt{WriteBatch} để ghi 14 documents (1 user + 1 profile + 1 preferences + 11 categories) trong 1 atomic operation.
  \item \textbf{Error mapping:} Firebase error codes (\texttt{email-already-in-use}, \texttt{weak-password}) được map sang thông báo tiếng Việt/Anh thân thiện qua \texttt{\_authErrorMessage()}.
\end{itemize}

\subsection{Chức năng 2: Trợ lý AI tài chính (AiFinanceService)}

Đây là chức năng đặc biệt nhất của ứng dụng. Hệ thống thu thập dữ liệu tài chính, xây dựng prompt, gọi Gemini API, và có cơ chế fallback khi AI không khả dụng.

\begin{lstlisting}[caption={AiFinanceService --- Xây dựng prompt và gọi Gemini API}]
static Future<String> generateAdvice({
  required List<MoneyTransaction> transactions,
  required List<Account> wallets,
}) async {
  // Fallback local khi khong co API key
  final fallback =
    _buildFallbackAdvice(transactions, wallets);
  if (!isConfigured) return fallback;

  // Thu thap du lieu tai chinh
  final totalBalance = wallets
    .fold<double>(0, (s, w) => s + w.balance);
  final monthTx = transactions
    .where((t) => !t.date.isBefore(monthStart));
  final income = monthTx
    .where((t) => t.isIncome)
    .fold<double>(0, (s, t) => s + t.amount);
  final expense = monthTx
    .where((t) => !t.isIncome)
    .fold<double>(0, (s, t) => s + t.amount);

  // Tim top 3 danh muc chi lon nhat
  final topExpenses = <String, double>{};
  for (final tx in monthTx.where((t) => !t.isIncome)) {
    topExpenses.update(
      tx.category, (v) => v + tx.amount,
      ifAbsent: () => tx.amount,
    );
  }

  // Xay dung prompt
  final prompt = StringBuffer()
    ..writeln('Ban la co van tai chinh ca nhan.')
    ..writeln('Tra loi bang tieng Viet, ngan gon.')
    ..writeln('Dua ra dung 3 goi y hanh dong.')
    ..writeln('')
    ..writeln('Du lieu nguoi dung:')
    ..writeln('- So du: $totalBalance VND')
    ..writeln('- Thu thang: $income VND')
    ..writeln('- Chi thang: $expense VND');

  // Goi API voi timeout 20s
  try {
    final response = await http.post(
      geminiUrl,
      headers: {'Content-Type': 'application/json'},
      body: jsonEncode({
        'contents': [
          {'parts': [{'text': prompt.toString()}]}
        ]
      }),
    ).timeout(const Duration(seconds: 20));

    // Parse response
    if (response.statusCode >= 200
        && response.statusCode < 300) {
      return parseGeminiText(response.body);
    }
    return fallback;
  } catch (_) {
    return fallback; // Luon tra ve ket qua
  }
}
\end{lstlisting}

\textbf{Giải thích kỹ thuật:}
\begin{itemize}
  \item \textbf{Strategy Pattern:} Hệ thống luôn có 2 chiến lược --- Gemini API hoặc phân tích local. Ứng dụng \textit{không bao giờ} để người dùng thấy trang trống.
  \item \textbf{Prompt Engineering:} Prompt được thiết kế cẩn thận: yêu cầu tiếng Việt, đúng 3 gợi ý, mỗi gợi ý tối đa 2 câu, ưu tiên hành động thực tế trong 7 ngày.
  \item \textbf{Data Privacy:} Chỉ gửi dữ liệu tổng hợp (tổng thu, tổng chi, top danh mục), không gửi chi tiết từng giao dịch.
  \item \textbf{Timeout 20 giây:} Đảm bảo UX tốt --- nếu API chậm, chuyển sang fallback thay vì để người dùng chờ.
  \item \textbf{Bảo mật API key:} Sử dụng \texttt{String.fromEnvironment()} + \texttt{--dart-define}, API key không xuất hiện trong source code.
\end{itemize}

\subsection{Chức năng 3: Đồng bộ dữ liệu tự động}

\texttt{MainShellPage} quản lý đồng bộ dữ liệu với cơ chế \textbf{periodic auto-sync}:

\begin{lstlisting}[caption={MainShellPage --- Auto-sync mechanism}]
class _MainShellPageState
    extends ConsumerState<MainShellPage> {
  Timer? _syncTimer;
  bool _syncInProgress = false;

  @override
  void initState() {
    super.initState();
    WidgetsBinding.instance.addPostFrameCallback((_) {
      _startAutoSync();
      _syncAllData(); // Sync ngay khi mo app
    });
  }

  void _startAutoSync() {
    _syncTimer?.cancel();
    _syncTimer = Timer.periodic(
      const Duration(seconds: 8), (_) {
        if (!mounted) return;
        _syncAllData();
      },
    );
  }

  Future<void> _syncAllData() async {
    if (_syncInProgress) return; // Tranh chong cheo
    _syncInProgress = true;
    try {
      // Dong thoi load accounts va transactions
      await Future.wait([
        AccountsRepository.instance
          .loadAccounts(forceRefresh: true),
        TransactionsRepository.instance
          .loadTransactions(forceRefresh: true),
      ]);

      if (!mounted) return;

      // Cap nhat Riverpod providers
      final accounts =
        AccountsRepository.instance.accounts;
      final transactions =
        TransactionsRepository.instance.transactions;

      ref.read(overviewWalletsProvider.notifier)
        .state = List.unmodifiable(accounts);
      ref.read(overviewTransactionsProvider.notifier)
        .state = List.unmodifiable(transactions);

      // Invalidate de cac widget khac rebuild
      ref.invalidate(walletsProvider);
      ref.invalidate(transactionsControllerProvider);
    } finally {
      _syncInProgress = false;
    }
  }

  @override
  void dispose() {
    _syncTimer?.cancel();
    super.dispose();
  }
}
\end{lstlisting}

\textbf{Giải thích kỹ thuật:}
\begin{itemize}
  \item \textbf{Tại sao polling thay vì realtime listener?} Realtime listeners (\texttt{snapshots()}) giữ kết nối WebSocket mở liên tục, tốn bandwidth và tăng Firestore billing. Polling mỗi 8 giây giảm 90\% reads so với realtime.
  \item \textbf{Mutex pattern:} Biến \texttt{\_syncInProgress} đảm bảo không có 2 sync chạy đồng thời.
  \item \textbf{Future.wait():} Load accounts và transactions song song, giảm 50\% thời gian chờ.
  \item \textbf{Unmodifiable lists:} Sử dụng \texttt{List.unmodifiable()} để state là immutable, đúng quy tắc Riverpod.
  \item \textbf{mounted check:} Kiểm tra widget còn tồn tại trước khi cập nhật state, tránh crash khi navigate đi.
\end{itemize}

\subsection{Chức năng 4: Repository Pattern (Tầng dữ liệu)}

\begin{lstlisting}[caption={TransactionsRepository --- Singleton + Cache + CRUD}]
class TransactionsRepository {
  // Singleton pattern
  static final instance = TransactionsRepository._();
  TransactionsRepository._();

  // In-memory cache
  List<MoneyTransaction> _cache = [];
  List<MoneyTransaction> get transactions => _cache;

  Future<void> loadTransactions({
    bool forceRefresh = false,
  }) async {
    if (!forceRefresh && _cache.isNotEmpty) return;

    final uid = AuthService.currentUserId;
    if (uid == null) return;

    final snapshot = await FirebaseFirestore.instance
      .collection('users')
      .doc(uid)
      .collection('transactions')
      .orderBy('date', descending: true)
      .get();

    _cache = snapshot.docs
      .map((doc) =>
        MoneyTransaction.fromFirestore(doc))
      .toList();
  }

  Future<void> addTransaction(
    MoneyTransaction tx,
  ) async {
    final uid = AuthService.currentUserId;
    await FirebaseFirestore.instance
      .collection('users')
      .doc(uid)
      .collection('transactions')
      .add(tx.toMap());
    _cache.insert(0, tx); // Cap nhat cache
  }
}
\end{lstlisting}

% ── 3.2 DEMO GIAO DIỆN ──────────────────────────────────────────────────────

\section{Demo Giao diện (UI Gallery)}

Dưới đây là ảnh chụp màn hình các trang chính, trình bày theo từng Usecase.

% Hướng dẫn: Thay placeholder bằng ảnh thực
% \includegraphics[width=0.4\textwidth]{images/screen_name.png}

\subsection{UC01--02: Splash Screen \& Đăng nhập}

\begin{figure}[H]
\centering
\fbox{\begin{minipage}{0.35\textwidth}
\centering\vspace{4cm}
\textit{Ảnh: Splash Screen\\với animation logo}
\vspace{4cm}
\end{minipage}}
\hspace{1cm}
\fbox{\begin{minipage}{0.35\textwidth}
\centering\vspace{4cm}
\textit{Ảnh: Login Page\\(đăng nhập / đăng ký)}
\vspace{4cm}
\end{minipage}}
\caption{Splash Screen có animation chuyên nghiệp và trang Đăng nhập hỗ trợ cả đăng nhập nhanh bằng sinh trắc học}
\label{fig:ui-auth}
\end{figure}

\subsection{UC05: Tổng quan tài chính (Dashboard)}

\begin{figure}[H]
\centering
\fbox{\begin{minipage}{0.35\textwidth}
\centering\vspace{4cm}
\textit{Ảnh: Overview Page\\tổng số dư, biểu đồ tháng}
\vspace{4cm}
\end{minipage}}
\hspace{1cm}
\fbox{\begin{minipage}{0.35\textwidth}
\centering\vspace{4cm}
\textit{Ảnh: Overview Page (cuộn xuống)\\gợi ý AI, giao dịch gần đây}
\vspace{4cm}
\end{minipage}}
\caption{Dashboard hiển thị thông tin tài chính tổng quan với biểu đồ cột thu/chi và gợi ý AI}
\label{fig:ui-overview}
\end{figure}

\subsection{UC03: Quản lý giao dịch}

\begin{figure}[H]
\centering
\fbox{\begin{minipage}{0.35\textwidth}
\centering\vspace{4cm}
\textit{Ảnh: Transaction Screen\\danh sách giao dịch}
\vspace{4cm}
\end{minipage}}
\hspace{1cm}
\fbox{\begin{minipage}{0.35\textwidth}
\centering\vspace{4cm}
\textit{Ảnh: Add Transaction Form\\form thêm giao dịch}
\vspace{4cm}
\end{minipage}}
\caption{Danh sách giao dịch với bộ lọc và form thêm giao dịch trực quan}
\label{fig:ui-transactions}
\end{figure}

\subsection{UC04: Quản lý ngân sách}

\begin{figure}[H]
\centering
\fbox{\begin{minipage}{0.35\textwidth}
\centering\vspace{4cm}
\textit{Ảnh: Budgets Page\\thanh tiến độ chi tiêu}
\vspace{4cm}
\end{minipage}}
\hspace{1cm}
\fbox{\begin{minipage}{0.35\textwidth}
\centering\vspace{4cm}
\textit{Ảnh: Budget Form\\form tạo ngân sách}
\vspace{4cm}
\end{minipage}}
\caption{Quản lý ngân sách với thanh tiến độ trực quan và cảnh báo vượt hạn mức}
\label{fig:ui-budgets}
\end{figure}

\subsection{UC06: Báo cáo thống kê}

\begin{figure}[H]
\centering
\fbox{\begin{minipage}{0.35\textwidth}
\centering\vspace{4cm}
\textit{Ảnh: Reports Screen\\biểu đồ tròn chi tiêu}
\vspace{4cm}
\end{minipage}}
\hspace{1cm}
\fbox{\begin{minipage}{0.35\textwidth}
\centering\vspace{4cm}
\textit{Ảnh: Reports Screen\\danh sách \% danh mục}
\vspace{4cm}
\end{minipage}}
\caption{Báo cáo chi tiết với biểu đồ tròn phân bổ chi tiêu và phần trăm từng danh mục}
\label{fig:ui-reports}
\end{figure}

\subsection{UC07--08: Cài đặt \& Quản lý ví/danh mục}

\begin{figure}[H]
\centering
\fbox{\begin{minipage}{0.35\textwidth}
\centering\vspace{4cm}
\textit{Ảnh: Profile/Settings Page}
\vspace{4cm}
\end{minipage}}
\hspace{1cm}
\fbox{\begin{minipage}{0.35\textwidth}
\centering\vspace{4cm}
\textit{Ảnh: Manage Wallets Page}
\vspace{4cm}
\end{minipage}}
\caption{Trang cài đặt và quản lý ví/tài khoản}
\label{fig:ui-settings}
\end{figure}

\subsection{UC09: Trợ lý AI \& Giới thiệu}

\begin{figure}[H]
\centering
\fbox{\begin{minipage}{0.35\textwidth}
\centering\vspace{4cm}
\textit{Ảnh: AI Chat Page\\chat với trợ lý AI}
\vspace{4cm}
\end{minipage}}
\hspace{1cm}
\fbox{\begin{minipage}{0.35\textwidth}
\centering\vspace{4cm}
\textit{Ảnh: About Page\\thông tin ứng dụng}
\vspace{4cm}
\end{minipage}}
\caption{Chat AI tài chính và trang giới thiệu ứng dụng}
\label{fig:ui-ai-about}
\end{figure}

% ── 3.3 ĐÁNH GIÁ RESPONSIVENESS ─────────────────────────────────────────────

\section{Đánh giá tính đáp ứng (Responsiveness)}

\subsection{Kỹ thuật responsive trong dự án}

\begin{table}[H]
\centering
\caption{Các kỹ thuật responsive được áp dụng}
\begin{tabularx}{\textwidth}{|l|X|}
\hline
\textbf{Kỹ thuật} & \textbf{Mô tả chi tiết} \\
\hline
\texttt{Flexible} \& \texttt{Expanded} & Tự động co giãn widget theo chiều rộng màn hình, đảm bảo layout không bị tràn \\
\hline
\texttt{CustomScrollView} + \texttt{Slivers} & Cuộn mượt, SliverAppBar co giãn từ expanded (240px) sang pinned (56px), hiệu ứng parallax \\
\hline
\texttt{SafeArea} & Tránh UI bị che bởi notch, status bar, home indicator trên các thiết bị mới \\
\hline
\texttt{MediaQuery} & Lấy kích thước màn hình, padding hệ thống, text scale factor \\
\hline
\texttt{LayoutBuilder} & Responsive layout dựa trên constraints thực tế của parent widget \\
\hline
\texttt{Wrap} widget & Tự động xuống dòng cho các chip, tag khi không đủ chiều rộng \\
\hline
\end{tabularx}
\end{table}

\subsection{Kết quả kiểm thử trên các thiết bị}

\begin{table}[H]
\centering
\caption{Bảng kiểm thử responsiveness chi tiết}
\label{tab:responsiveness}
\begin{tabularx}{\textwidth}{|l|l|l|c|X|}
\hline
\textbf{Thiết bị} & \textbf{Màn hình} & \textbf{OS} & \textbf{OK?} & \textbf{Ghi chú} \\
\hline
Samsung Galaxy A54 & 6.4" FHD+ & Android 14 & ✓ & Thiết bị test chính \\
\hline
iPhone 15 (Simulator) & 6.1" Super Retina & iOS 17 & ✓ & Kiểm tra notch, safe area \\
\hline
Pixel 7 (Emulator) & 6.3" FHD+ & Android 14 & ✓ & Tham chiếu Google design \\
\hline
Galaxy S21 FE & 6.4" FHD+ & Android 13 & ✓ & Thiết bị đời cũ hơn \\
\hline
iPad Air (Simulator) & 10.9" & iPadOS 17 & ✓ & Layout vẫn đẹp trên tablet \\
\hline
Chrome (Desktop) & 1920$\times$1080 & Windows 11 & ✓ & Web version \\
\hline
Edge (Desktop) & 1366$\times$768 & Windows 11 & ✓ & Màn hình laptop nhỏ \\
\hline
\end{tabularx}
\end{table}

\subsection{Hỗ trợ Dark Mode}

Ứng dụng hỗ trợ đầy đủ 2 chế độ giao diện:

\begin{figure}[H]
\centering
\fbox{\begin{minipage}{0.35\textwidth}
\centering\vspace{3cm}
\textit{Ảnh: Light Mode}
\vspace{3cm}
\end{minipage}}
\hspace{1cm}
\fbox{\begin{minipage}{0.35\textwidth}
\centering\vspace{3cm}
\textit{Ảnh: Dark Mode}
\vspace{3cm}
\end{minipage}}
\caption{So sánh giao diện Light Mode và Dark Mode}
\label{fig:light-dark}
\end{figure}

Triển khai Dark Mode sử dụng \texttt{ColorScheme.fromSeed()} với cùng seed color nhưng brightness khác nhau:

\begin{lstlisting}[caption={Theme configuration trong app.dart}]
// Light theme
theme: ThemeData(
  colorScheme: ColorScheme.fromSeed(
    seedColor: const Color(0xFF0D7377),
    secondary: const Color(0xFF6C63FF),
    tertiary: const Color(0xFF14A085),
    brightness: Brightness.light,
  ),
  useMaterial3: true,
),
// Dark theme (cung seed, khac brightness)
darkTheme: ThemeData(
  colorScheme: ColorScheme.fromSeed(
    seedColor: const Color(0xFF0D7377),
    secondary: const Color(0xFF6C63FF),
    tertiary: const Color(0xFF14A085),
    brightness: Brightness.dark,
  ),
  useMaterial3: true,
),
\end{lstlisting}

Toàn bộ UI sử dụng \texttt{Theme.of(context).colorScheme} thay vì hardcode màu, đảm bảo tự động chuyển đổi khi thay đổi theme.
